\begin{table}[H]
\centering
\tiny
\renewcommand{\arraystretch}{1.15}
\setlength{\tabcolsep}{3pt}
\resizebox{\textwidth}{!}{%
\begin{tabular}{|l|r|r|l|r|}
\hline
\multicolumn{3}{|c|}{\textbf{Desarrollo de Servicio}} & \multicolumn{2}{c|}{\textbf{Periodo de Recuperación}} \\
\multicolumn{3}{|c|}{\textbf{Valor Actual Neto}} & \multicolumn{2}{c|}{\textbf{con los Flujos Descontados}} \\ \hline
 & \textbf{FLUJOS NETOS} & \textbf{FLUJOS DESCONTADOS} &  &  \\ \hline
INVERSIÓN INICIAL & -19,154.29 & -19,154.29 & \textbf{Inversión} & \textbf{19,154.29} \\ \hline
FLUJO 1ER AÑO & -70,980.87 & -62,157.63 &  & -70,980.87 \\ \hline
FLUJO 2DO AÑO & 6,310.24 & 4,838.98 &  & -64,670.63 \\ \hline
FLUJO 3ER AÑO & 146,441.88 & 98,338.75 &  & 81,771.26 \\ \hline
FLUJO 4TO AÑO & 78,368.84 & 46,084.61 &  & 160,140.10 \\ \hline
FLUJO 5TO AÑO & 360,022.75 & 185,393.97 &  & 520,162.85 \\ \hline
FLUJO 6TO AÑO & 533,659.88 & 240,648.60 &  & 1,053,922.73 \\ \hline
 &  &  &  &  \\ \hline
\textbf{Criterio de decisión} &  & \textbf{VAN =} & \textbf{S/} & \textbf{493,992.97} \\ \hline
\multicolumn{5}{|l|}{\textbf{RESULTADO, Nuestra Van es positiva y mayor que cero.}} \\ \hline
\multicolumn{5}{|l|}{\textbf{PROYECTO SE ACEPTA}} \\ \hline
\end{tabular}%
}

\vspace{0.2cm}
\textbf{Figura 10.29.} Valor Neto Actual y Periodo de Recuperación \textbf{[Elaboración Propia]}
\end{table}

\textbf{Tasa Interna de Retorno}

La tasa interna de retorno se calcula en Excel así =TIR (celda de la inversión en negativo: último flujo) =IRR(A1:A5)

\begin{table}[H]
\centering
\scriptsize
\renewcommand{\arraystretch}{1.2}
\setlength{\tabcolsep}{6pt}
\begin{tabular}{|p{8cm}|p{4cm}|}
\hline
\multicolumn{2}{|c|}{\textbf{Desarrollo de Servicio}} \\
\multicolumn{2}{|c|}{\textbf{Tasa Interna de Retorno}} \\ \hline
\textbf{Criterio de decisión TIR} & \textbf{TIR= \hfill 88.56\%} \\ \hline
RESULTADO \hfill \textbf{k<TIR} &  \\ \hline
PROYECTO SE ACEPTA &  \\ \hline
\end{tabular}

\vspace{0.2cm}
\textbf{Figura 10.30.} Tasa interna de retorno \textbf{[Elaboración Propia]}
\end{table}

\textbf{Índice de Rentabilidad}

El Índice de Rentabilidad es la suma de los flujos descontados entre el valor de la inversión:

\begin{table}[H]
\centering
\scriptsize
\renewcommand{\arraystretch}{1.2}
\setlength{\tabcolsep}{6pt}
\begin{tabular}{|p{8cm}|p{4cm}|}
\hline
\multicolumn{2}{|c|}{\textbf{Índice de Rentabilidad}} \\ \hline
\textbf{Criterio de decisión IR} & \textbf{IR= \hfill 26.79} \\ \hline
RESULTADO \hfill \textbf{IR > 1} &  \\ \hline
PROYECTO SE ACEPTA &  \\ \hline
\end{tabular}

\vspace{0.2cm}
\textbf{Figura 10.31.} Índice de Rentabilidad \textbf{[Elaboración Propia]}
\end{table}